\documentclass[instructions.tex]{subfiles}
\begin{document}

\chapter{Il est plus difficile de se corriger des habitudes contractées par le penchant.}


En suivant ses penchans vicieux, on fait à son âme des plaies qui sont bien plus incurables, soit du côté du cœur, soit du côté de l'esprit.

\primo Du côté du cœur, il est plus difficile de s'en repentir. On ne se repent pas aisément d'un péché qui vient d'un penchant qu'on aime, ou qu'on ne veut point connaître. On est fâché, dit-on, d'avoir péché; on le dit, mais souvent on se trompe. On n'est pas toujours fâché d'avoir fait le péché, mais on est fâché de ce qu'il y a du péché dans ce qu'on a fait, parce qu'on voudrait qu'il n'y eût point de péché dans ce qui plaît. On est fâché, parce qu'il faut s'en confesser et s'en corriger; mais en est-on fâché parce qu'il offense Dieu? C'est sur ce point qu'il faudrait sonder son cœur.

Nous disons que le péché nous déplaît; et à qui ne déplaît-il pas? il déplaît même aux plus libertins. On parlerait avec plus de vérité, si l'on disait que c'est la gêne qui nous déplaît, et que la violence qu'il faut se faire pour se mettre au-dessus du qu'en dira-t-on et pour se corriger, nous déplaît encore plus que nos péchés.

On craint bien plus d'être raillé du monde, de faire quelques efforts et de s'incommoder, qu'on ne craint de pécher. Il est donc vrai qu'il est difficile de quitter un vice, surtout quand il est d'habitude, lorsque le penchant a gagné le cœur.

\secundo Il n'est pas moins difficile du côté de l'esprit; car comment combattre un penchant qui aveugle la raison, qui étouffe les remords, et qu'on ne veut point connaître?

La conscience, les lumières de la foi, la parole de Dieu, ne dissipent guère les illusions et le charme qui séduisent un esprit prévenu par la passion et par le penchant. Les instructions pourraient-elles désabuser une personne qui croit le contraire de ce qu'on lui dit? A quoi sert de présenter de la lumière à un homme qui ferme les yeux? et comment guérir celui qui croit n'être pas malade?

N'est-ce pas pour cette raison que les jeunes gens, dont le penchant est si porté au plaisir et à la dissipation, ouvrent si difficilement les yeux sur les fréquentations et sur les attaches qui les perdent? que l'impudique, loin de comprendre l'horreur de ses crimes honteux, les traite de bagatelles ou de faiblesses? que plusieurs personnes de qualité, passionnées pour le plaisir et pour la gloire, loin d'ouvrir les yeux sur les dangers de leur condition, traitent de bienséance leurs dépenses fastueuses, leur vie molle et leur ambition? que les militaires regardent sans horreur le crime détestable du duel comme un point d'honneur, et ne voient pas qu'il est l'effet d'une fureur diabolique et d'un orgueil insensé qui les aveuglent, et que Dieu réprouve?

N'est-ce pas aussi pour ne pas connaître leurs penchans déréglés, que tant de personnes attachées à leurs biens, emportées quand elles perdent, dures envers le pauvre et le débiteur, ne voient pas que l'avarice et l'esprit d'intérêt les possèdent? que tant de gens portés à la médisance, à la satire, qui s'offensent de tout, qui refusent de se réconcilier et de voir ceux qui les ont offensés, ne s'aperçoivent pas qu'ils sont sans charité; et, vivant, sans le savoir, dans une jalousie et une rancune habituelles, se persuadent, malgré le scandale qu'ils donnent, qu'ils sont innocens?

N'est-ce pas pour la même raison que tant de joueurs, de fainéans, d'intempérans qui fréquentent habituellement les tavernes, se persuadent qu'ils ne font point de mal? que tant de gens avides regardent comme innocens le trafic et l'exercice du commerce dans les saints jours, et se croient permis les monopoles, les usures, les supercheries et les chicanes; que tant de gens, impatients quand on les contredit, pleins de hauteur dans leurs sentimens, qui prennent feu pour une parole sur le compte de leur famille, ne voient pas que c'est un orgueil secret qui les domine? N'est-ce pas pour cette raison que tant de pères et de mères aveuglés par l'esprit du monde ou par un amour excessif de leurs enfans, souffrent dans leurs familles des désordres, et des intrigues qu'ils ne veulent ni savoir, ni croire, ni empêcher? N'est-ce pas enfin parce qu'on ne connaît pas ses penchans secrets, que tant de gens s'accusent si mal de leurs défauts essentiels, pendant qu'ils s'accusent si scrupuleusement de leurs fautes légères?

Le prophète avait donc bien raison de dire : Qui est celui qui connaît ses péchés? L'amour-propre, le penchant pour les plaisirs et les choses de la terre, sont comme un nuage épais qui obscurcit la raison, sont comme un souffle venimeux qui peu à peu éteint les lumières de la foi. Certains péchés grossiers, pour lesquels on n'a pas de penchans, se font sentir : on les connaît, on s'en corrige; mais les péchés que le penchant fomente, sont ordinairement les plus dangereux pour le salut : on n'en connaît souvent ni la cause, ni la malice, ni les remèdes, parce qu'on ne veut pas les connaître\footnote{Delicta quis intelligit? Ps. 18.}. Veillez donc sur vous-même, sondez votre cœur, instruisez-vous et laissez-vous instruire.


\section{Résolutions}

\primo Quelque difficile qu'il soit de rompre les habitudes pernicieuses appuyées sur de mauvais penchans, je suis résolu de travailler dès aujourd'hui à détruire celles qui règnent en moi.... Quelles sont-elles?

\secundo J'aurai recours souvent, à cet effet, à la prière, aux sacremens, à la méditation des grandes vérités, et je veillerai continuellement sur moi à l'égard de ces funestes habitudes.

Venez à mon secours, ô mon Dieu! afin que je parvienne à tarir à jamais ces sources fécondes de transgressions.

\end{document}
